% Part: normal-modal-logic
% Chapter: axioms-systems
% Section: consistency

\documentclass[../../../include/open-logic-section]{subfiles}

\begin{document}

\olfileid{nml}{prf}{con}

\olsection{الاتساق}

الاتساق خاصية مهمة لمجموعات ال!!{formula}s. وتكون مجموعة من
ال!!{formula}s غير متسقة إذا كان من الممكن أن !!{derivable} منها تناقض،
مثل~$\lfalse$؛ وتكون متسقة فيما عدا ذلك. وإذا كانت مجموعة غير متسقة،
فلا يمكن أن تكون !!{formula}sها كلها صادقة في نموذج عند عالم واحد.
وسنثبت في مبرهنة التمام العكس: كل مجموعة متسقة صادقة عند عالم في نموذج،
هو تحديدًا ``النموذج القانوني''.

\begin{defn}
  تكون المجموعة~$\Gamma$ \emph{متسقة} بالنسبة إلى النسق~$\Sigma$، أو
  كما سنقول، \emph{$\Sigma$-متسقة}، إذا وفقط إذا كان
  $\Gamma \Proves/[\Sigma] \lfalse$.
\end{defn}

فمثلًا، المجموعة~$\{\Box(p \lif q), \Box p, \lnot\Box q\}$ متسقة
بالنسبة إلى منطق القضايا، لكنها ليست~\Log{K}-متسقة. وبالمثل، ليست
المجموعة~$\{\Diamond p, \Box\Diamond p \lif q, \lnot q\}$
\Log{K5}-متسقة.

\begin{prop}\ollabel{prop:consistencyfacts}
  لتكن~$\Gamma$ مجموعة من ال!!{formula}s. عندئذ:
  \begin{enumerate}
  \item تكون~$\Gamma$ $\Sigma$-متسقة إذا وفقط إذا وجدت
    !!{formula}~$!A$ بحيث~$\Gamma \Proves/[\Sigma] !A$.
  \item \ollabel{prop:consistencyfacts-b}%
    $\Gamma \Proves[\Sigma] !A$ إذا وفقط إذا لم تكن
    $\Gamma \cup \{\lnot!A\}$ $\Sigma$-متسقة.
  \item \ollabel{prop:consistencyfacts-c}%
    إذا كانت~$\Gamma$ $\Sigma$-متسقة، فلكل !!{formula}~$!A$ تكون إما
    $\Gamma \cup \{!A\}$ $\Sigma$-متسقة أو
    $\Gamma \cup \{\lnot!A\}$ $\Sigma$-متسقة.
  \end{enumerate}
\end{prop}

\begin{proof}
  تنتج هذه الوقائع بسهولة من منطق القضايا الكلاسيكي. ونورد حجة
  \olref{prop:consistencyfacts-c}. نستدل بالمقابل، فنفترض أن لا
  $\Gamma \cup \{!A\}$ ولا~$\Gamma \cup \{\lnot!A\}$
  $\Sigma$-متسقة. عندئذ، بحسب~\olref{prop:consistencyfacts-b}، لدينا
  $\Gamma, !A \Proves[\Sigma] \lfalse$ و
  $\Gamma, \lnot !A \Proves[\Sigma] \lfalse$. وبمبرهنة الاستنتاج،
  $\Gamma \Proves[\Sigma] !A \lif \lfalse$ و
  $\Gamma \Proves[\Sigma] \lnot!A \lif \lfalse$. لكن
  $(!A \lif \lfalse) \lif ((\lnot!A \lif \lfalse) \lif \lfalse)$
  مثال إحلال لقضية صادقة صدقًا تحليليًا؛ ومن ثم، بحسب
  \olref[prp]{prop:derivabilityfacts}\olref[prp]{prop:derivabilityfacts-ruleT}،
  $\Gamma \Proves[\Sigma] \lfalse$.
\end{proof}

\end{document}
